\chapter{Networking: Services, DNS, Ingress}
\label{ch:networking}

\section{The problem pods create}

Pods have cluster-routable IPs --- and lose them at every replacement.
Chapter~15 made pods disposable; something stable must sit in front.
That something is the \emph{Service}: a fixed virtual IP and DNS name,
backed by a live list of pod endpoints selected by label:

\begin{lstlisting}[language=yaml]
apiVersion: v1
kind: Service
metadata: { name: course-service }
spec:
  selector: { app: course-service }     (*@\co{1}@*)
  ports:
    - { port: 8082, targetPort: 8082 }  (*@\co{2}@*)
\end{lstlisting}

\co{1}~The endpoint controller continuously lists \emph{Ready} pods
matching this selector --- readiness (Chapter~17) is the membership
test. \co{2}~Traffic to the Service's port is forwarded to a chosen
endpoint's targetPort.

Cluster DNS gives every Service a name: \code{course-service} within
the namespace, \code{course-service.ts.svc.cluster.local} from
anywhere. This is why the Kubernetes env overrides look so much like
the Compose ones --- both platforms resolve a service's name on a
private network; the project deliberately named its k8s Services after
its Compose services so the overrides match.

The default Service type, \code{ClusterIP}, is reachable \emph{only
inside} the cluster --- a security property Part~II leaned on for the
password-less Mongo. \code{NodePort} and \code{LoadBalancer} exist to
expose Services directly; this project instead funnels everything
through one Ingress.

\section{Ingress and Traefik}

A Service exposes one thing; an \emph{Ingress} is the L7 front door:
host- and path-based HTTP routing, later TLS termination, declared as
data and executed by an \emph{Ingress controller} --- in k3s, the
bundled Traefik watching Ingress objects. From
\code{api-gateway.yaml}:

\begin{lstlisting}[language=yaml]
kind: Ingress
spec:
  rules:
    - host: ts.local                    (*@\co{1}@*)
      http:
        paths:
          - path: /api                  (*@\co{2}@*)
            pathType: Prefix
            backend:
              service:
                name: api-gateway
                port: { number: 8080 }
          # /oauth2 and /login/oauth2 likewise -> api-gateway
\end{lstlisting}

\begin{description}
  \item[\co{1}] Host-based rule: this applies to requests whose
  \code{Host} header says \code{ts.local} --- a name that exists only in
  the Windows hosts file, pointing at the server's tailnet IP. Multiple
  hosts can share the cluster's single 80/443.
  \item[\co{2}] Path fan-out. Only three prefixes enter, all to the
  gateway; \code{/} is left free for the frontend's own Ingress
  (Chapter~26), which Traefik merges into the same host's routing table.
\end{description}

\section{The full path of one request}

\begin{figure}[h]
\centering
\begin{tikzpicture}[node distance=3.5mm and 5.5mm]
  \node[boxdark] (b) {Browser};
  \node[box, right=of b] (t) {Traefik};
  \node[box, right=of t] (gws) {Service\\\scriptsize api-gateway};
  \node[box, right=of gws] (gw) {gateway pod};
  \node[box, below=of gw] (cs) {Service\\\scriptsize course-service};
  \node[box, left=of cs] (csp) {course pod};
  \draw[flow] (b) -- node[lbl,above]{\scriptsize Host: ts.local} (t);
  \draw[flow] (t) -- node[lbl,above]{\scriptsize /api} (gws);
  \draw[flow] (gws) -- (gw);
  \draw[flow] (gw) -- node[lbl,right]{\scriptsize lb:// via Eureka}(cs);
  \draw[flow] (cs) -- (csp);
\end{tikzpicture}
\caption{Two routers, two discovery systems, one request.}
\end{figure}

Worth pausing on: the gateway's Eureka lookup returns a \emph{pod IP}
(Chapter~3's prefer-ip-address), so hop four technically bypasses the
course-service Service. Both discovery systems run in parallel ---
Eureka for gateway-to-service, Service DNS for everything else
(config-server, databases, Kafka). Chapter~27's simplification
collapses them into one.

\begin{pitfall}[the Service that selects nothing]
Symptom: \code{curl} to a Service name times out; pods are Running and
healthy. Cause \#1: label mismatch --- the Service selector and the pod
template labels differ by a typo; the endpoint list is empty and
nothing complains. Check: \code{kubectl -n ts get endpoints
course-service} --- an empty ENDPOINTS column is this bug. Cause \#2:
endpoints exist but the pods are not Ready --- same command shows them
under \code{NotReadyAddresses}. The Service abstraction fails
\emph{silently}; the endpoints object is where it becomes visible.
\end{pitfall}

\begin{decision}[one Ingress host + hosts-file, not NodePorts]
NodePorts (one weird port per service) would work over Tailscale
without any DNS games --- and produce URLs like
\code{100.85.66.43:30082} that leak topology, dodge the gateway, and
can't do TLS sanely. One Ingress host means one address to secure and
a browser experience identical to production; the \code{ts.local}
hosts-file entry is the only client-side cost. The upgrade path ---
a real domain plus cert-manager, or a Cloudflare Tunnel --- changes the
Ingress host and nothing else.
\end{decision}

\section{Outside the project}

On EKS the same Ingress object, annotated for the AWS controller,
provisions an actual ALB; on OpenShift it becomes a \code{Route} ---
same idea, different noun. The piece this cluster omits:
\textbf{NetworkPolicies} --- firewall rules for pod-to-pod traffic
(``only the gateway may reach course-service:8082''), which is the
in-cluster enforcement the trusted-edge pattern of Chapter~5 called
its remaining soft spot. Flat-network trust is normal for a homelab
and negligent for a bank.

\begin{quiz}
\begin{enumerate}
  \item Pod IPs change at every replacement, yet
  \code{course-service:8082} always works. Name each mechanism between
  the name and the pod's socket.
  \item How does readiness interact with Service endpoints, and what
  does that mean during a rolling update?
  \item Diagnose: Service exists, pods Ready, \code{get endpoints}
  shows an empty list. What is wrong, and why did nothing error?
  \item Why does \code{ts.local} require a hosts-file entry, and what
  replaces that entry in the two production upgrade paths?
\end{enumerate}
\end{quiz}
